% =====================================================================
%  NSIGII Protocol Specification
%  OBINEXUS-NSIGII-SPEC  v0.1 (Draft)
%  Derived from the NSIGII transcript corpus (6 sessions, Jan-Jun 2026)
%  Companion document: NSIGII-SPEC.md
%  Build:  pdflatex nsigii-spec.tex  (run twice for the ToC)
% =====================================================================
\documentclass[11pt,a4paper]{article}

\usepackage[T1]{fontenc}
\usepackage[utf8]{inputenc}
%\usepackage{lmodern}
\usepackage[margin=1in]{geometry}
\usepackage{amsmath,amssymb}
\usepackage{booktabs}
\usepackage{longtable}
\usepackage{array}
\usepackage{enumitem}
\usepackage{xcolor}
\usepackage{fancyvrb}
\usepackage{tcolorbox}
\usepackage[hidelinks]{hyperref}
\usepackage{textcomp}

\definecolor{errcol}{HTML}{8A3324}
\definecolor{rulecol}{HTML}{1F3A5F}

\newcommand{\MUST}{\textsc{must}}
\newcommand{\MUSTNOT}{\textsc{must not}}
\newcommand{\SHOULD}{\textsc{should}}
\newcommand{\MAY}{\textsc{may}}
\newcommand{\NSIGII}{\textsc{nsigii}}
\newcommand{\CISCO}{\textsc{cisco}}
\newcommand{\LAPIS}{\textsc{lapis}}
\newcommand{\OBI}{\texttt{OBINEXUS}}

\newtcolorbox{errata}[1]{colback=errcol!5,colframe=errcol!70,
  fonttitle=\bfseries,title={#1},boxrule=0.5pt,left=4pt,right=4pt,
  top=3pt,bottom=3pt}
\newtcolorbox{ruleblock}{colback=rulecol!5,colframe=rulecol!60,
  boxrule=0.5pt,left=4pt,right=4pt,top=3pt,bottom=3pt}

\newcommand{\seg}[1]{\textsc{#1}}

\title{\bfseries \NSIGII{} Protocol Specification\\[4pt]
\large Trident Packet Verification via Bipolar \CISCO{}\\
\large Order/Chaos Separation of Concerns}
\author{Nnamdi Michael Okpala \\ OBINexus Computing}
\date{Version 0.3 --- Draft \\ Compiled 25 August 2026}

\begin{document}
\maketitle

\begin{center}
\begin{tabular}{@{}ll@{}}
\toprule
Document ID & \texttt{OBINEXUS-NSIGII-SPEC} \\
Version & 0.3 --- Draft \\
Status & Working draft, derived from primary transcript corpus \\
Source & \texttt{NSIGII.txt} --- 6 recorded sessions, Jan--Jun 2026 \\
Companion & \texttt{NSIGII-SPEC.md} (Markdown, same content) \\
\bottomrule
\end{tabular}
\end{center}

\tableofcontents
\newpage

% =====================================================================
\section{Document Status and Provenance}

\subsection{What this document is}

This is a \emph{derived} specification. It is not a transcript and not an
independent design. Every normative statement traces to a claim made in the
\NSIGII{} recording corpus. Where the corpus is internally inconsistent, this
document states the inconsistency rather than silently choosing a side.

\subsection{Source corpus}

\begin{center}
\begin{tabular}{@{}llp{8.4cm}@{}}
\toprule
\textbf{Seg} & \textbf{Date in recording} & \textbf{Working title} \\
\midrule
S1 & 30 Jan 2026, 05:39 & Trident packet verification via bipolar \CISCO{} --- order/chaos separation of concerns \\
S2 & 13 Jan 2026, 10:15 & Filter/flash interdependency --- \CISCO{} sequence and series instruction execution order via inverse relay trident messages \\
S3 & 13 Jan 2026, 18:08 & \NSIGII{} overview --- human rights protocol, pointer problem, \textsc{ox-star} \\
S4 & 7 Feb 2026 + 2 Mar 2026 & Formalising suffering; here-and-now/stillness; radio jamming and the maybe lattice \\
S5 & undated + 12 Feb 2026 & Trilateral governance; reference implementation demonstration \\
S6 & 28 Jun 2026 + undated & Dimensional principality and non-linearity; the \NSIGII{} codec and language classifiers \\
\bottomrule
\end{tabular}
\end{center}

\subsection{Transcription caveat}

The corpus is machine-transcribed speech. Proper nouns, symbols and numbers are
frequently corrupted: \texttt{in sigi}, \texttt{insigi}, \texttt{insidity},
\texttt{ciggy} all denote \NSIGII{}; \texttt{Cisco} denotes \CISCO{} as defined
in \S\ref{sec:cisco}, not the network vendor; \texttt{rift} denotes the
\textsc{rift} toolchain; \texttt{lapis} denotes \LAPIS{}
(\S\ref{sec:lapis}). Numeric claims are reproduced in
Appendix~\ref{app:errata} with corrected values.

\subsection{Conformance language}

\MUST{}, \MUSTNOT{}, \SHOULD{}, \MAY{} are used in the RFC~2119 sense.
Statements marked \emph{(unresolved)} are recorded intent that is not yet
specifiable.

% =====================================================================
\section{Overview}

\subsection{Name}

\NSIGII{} --- spelled in the corpus as November--Sierra--India--Golf--India--India.
The name is drawn from Igbo (Niger--Congo). Igbo is described in the corpus as a
tonal, verb--noun--constructed language, and that construction is load-bearing:
the protocol's instruction model is a verb--noun pairing (\S\ref{sec:verbnoun}),
and its tonal distinction is the canonical example of a message whose meaning
depends on a channel the encoding must preserve.

The corpus gives the tonal example directly: \emph{eze} under one tone contour
means \textbf{king}; under another it means \textbf{teeth}. The protocol's
stated position is that ``the king has teeth'' --- the governance layer and the
enforcement layer are the same word under different tone. A codec that drops
tone drops meaning; this is the motivating case for schema-level rather than
message-level verification (\S\ref{sec:repair}).

\subsection{Purpose}

\NSIGII{} is a \textbf{command-and-control protocol for safety-critical delivery
of human-rights entitlements} --- stated in the corpus, repeatedly and without
metaphor, as \emph{food, water and shelter, here and now}. Its technical content
is in service of that: a message saying ``I need food at this address now'' must
survive corruption, must be verifiable by a third party, must not be silently
dropped, and must leave an auditable record that a duty-holder received it.

The protocol is explicitly \emph{not} an AI system. Natural-language
understanding here is structural (\S\ref{sec:lang}), not learned.

\subsection{Governance invariants --- the Three Nevers}
\label{sec:nevers}

Implementations \MUST{} satisfy three invariants (S5). They are \textbf{not
three independent rules}: each is the invariant of one \textbf{addressing mode}
(\S\ref{sec:addressing}). The Three Nevers and the addressing triad are the same
structure seen twice (author ratification, 25 August 2026).

\begin{ruleblock}
\begin{center}
\textbf{here and now --- never a toy}\\
\textbf{there and then --- never a weapon}\\
\textbf{where and whenever --- never a problem}
\end{center}
\end{ruleblock}

\begin{center}
\begin{tabular}{@{}>{\raggedright\arraybackslash}p{3.1cm}>{\raggedright\arraybackslash}p{2.6cm}>{\raggedright\arraybackslash}p{8.4cm}@{}}
\toprule
\textbf{Mode} & \textbf{Never} & \textbf{Requirement} \\
\midrule
\textsc{here and now} & never a toy & The system \MUSTNOT{} be shipped as a
  demonstration or a game. \\[2pt]
\textsc{there and then} & never a weapon & The system \MUSTNOT{} be usable as an
  instrument of harm against the person it serves. \\[2pt]
\textsc{where and whenever} & never a problem & The system \MUSTNOT{} crash and
  \MUSTNOT{} accumulate debt. It fails \emph{gracefully}: times out, resolves
  outstanding state, saves, closes. \\
\bottomrule
\end{tabular}
\end{center}

\paragraph{Why each never belongs to its mode.}

\begin{enumerate}[leftmargin=*]
  \item \textsc{here and now} $\rightarrow$ \textbf{never a toy.} The deictic
    mode is where the need is \emph{actual} --- food, at this address, now. Both
    dimensions bound means both are real. A toy is a thing played with in the
    present; the mode that owns the present is therefore the mode that must
    forbid play. The corpus: \emph{``it's not a toy, that's not how to play with
    it --- it's something that gives people food and shelter.''}
  \item \textsc{there and then} $\rightarrow$ \textbf{never a weapon.}
    \textbf{Displacement is what makes a weapon possible.} Targeting \emph{is} a
    there-and-then operation: a specified elsewhere, a specified later, acted on
    at a distance. Every reach the protocol has --- \textsc{ox-star} injection,
    guiding a craft in flight, disarming a running system
    (\S\ref{sec:oxstar}) --- lives in this mode. The prohibition is placed
    exactly where the reach is. A protocol that can act at a distance \MUSTNOT{}
    act at a distance to harm.
  \item \textsc{where and whenever} $\rightarrow$ \textbf{never a problem.} The
    unbound mode \emph{cannot fail into a fault, because it has no slot to
    miss.} A claim quantified over all place and time is never overdue, so it
    accrues no debt; it is not failing, it is unresolved. Graceful failure is not
    a behaviour bolted onto this mode --- it is what unboundedness \emph{is}.
    This is the same result \S\ref{sec:stillness} reaches from the other
    direction.
\end{enumerate}

\begin{errata}{Naming note (v0.3)}
v0.1--0.2 of this document rendered the first invariant as ``never fails.'' That
is the corpus's \emph{gloss}, not its name. S5 states \textbf{``never a
problem''} and then explains it as \emph{``it never fails --- it fails
gracefully, it times out, it resolves all issues and saves.''} The canonical name
is restored here.
\end{errata}

\begin{errata}{Tension (now largely resolved)}
S5 states ``never a weapon --- you cannot sue people'' while S5b develops
\NSIGII{} explicitly as a \emph{litigation system} (\S\ref{sec:litigation}). The
mode pairing gives the resolution a precise locus: ``never a weapon'' is the
invariant of \textsc{there and then} --- the protocol's own reach across space
and time. A claimant using their own verified record to enforce their own
entitlement is not the protocol reaching out to harm; it is the claimant acting
in \textsc{here and now}, under a different invariant entirely. See
\S\ref{sec:reconcile}.
\end{errata}

\subsection{Governance topology}

\NSIGII{} is \textbf{trilateral}, over a lattice of three governance modes:
\textbf{unilateral} (one party acts, the protocol records); \textbf{bilateral}
(two parties, the yes/no axis); \textbf{trilateral} (three parties, the
yes/no/\emph{maybe} axis, which is the protocol's native consensus form,
\S\ref{sec:rgb}).

The transport topology is \textbf{peer-to-peer}, not blockchain. The corpus
rejects blockchain explicitly while retaining the property it wants from it ---
that a record, once made, has teeth.

\subsection{Position in the OBINexus toolchain}

\begin{Verbatim}[fontsize=\small,frame=single]
riftlang.exe  ->  .so.a  ->  rift.exe  ->  gosilang
                    ^
                 nlink  ->  polybuild
\end{Verbatim}

\begin{itemize}[leftmargin=*]
  \item \textbf{RIFT} supplies the token model \NSIGII{} packets carry: the
    triplet \emph{(token type, token value, token data)} ---
    \S\ref{sec:triplet}.
  \item \texttt{.rift} files carry the macro/meta control sequences \NSIGII{}
    relays.
  \item \textbf{gosilang} is the stated implementation target for the
    clause/\texttt{then} control-flow model (\S\ref{sec:then}).
  \item The reference implementation demonstrated in S5b is C, built with
    \texttt{make}, with a Go codec (\S\ref{sec:codec}).
\end{itemize}

% =====================================================================
\section{Foundational Model}

\subsection{The six operators}
\label{sec:ops}

All state change is organised into three complementary pairs:

\begin{center}
\begin{tabular}{@{}lll@{}}
\toprule
\textbf{Pair} & \textbf{Constructive} & \textbf{Destructive} \\
\midrule
Existence   & create & destroy \\
Structure   & build  & break \\
Restoration & repair & renew \\
\bottomrule
\end{tabular}
\end{center}

\begin{ruleblock}
\emph{If you can build something you can break it; if you can create something
you can destroy it; if you can repair something you can renew it.}
\end{ruleblock}

Capability is symmetric. An implementation offering one half of a pair \MUST{}
account for the other half --- either by implementing it or by declaring it
structurally impossible (\S\ref{sec:enzyme}).

Note that \emph{repair} and \emph{renew} are both constructive in ordinary
usage. The corpus treats \textbf{renew} as the operator replacing a value the
receiver cannot resolve (an $\epsilon$ or $\theta$, \S\ref{sec:corrupt}), and
\textbf{repair} as the operator correcting a value the receiver resolved
wrongly. That distinction is preserved.

\subsection{Order and Chaos sequences}
\label{sec:orderchaos}

Two opposed instruction sequences run over the six operators.

\begin{align*}
\textsc{order} \;&:\; \text{create} \rightarrow \text{build} \rightarrow \text{repair} \rightarrow \text{renew} \\
\textsc{chaos} \;&:\; \text{destroy} \rightarrow \text{break} \rightarrow \text{renew} \rightarrow \text{repair}
\end{align*}

\begin{ruleblock}
\textbf{\textsc{chaos} is not an error state.} Chaos is the \emph{observable};
order is inferred from it. \emph{``You check chaos for order just by observing
the dimension of order and chaos where the program was running.''}
\end{ruleblock}

Every channel carries a status of \texttt{ORDER} or \texttt{CHAOS}. The S5b
demonstration shows a receiver reporting status \texttt{chaos} and resolving
toward \texttt{order}: nominal behaviour, not a fault.

\subsection{The isolation principle (enzyme model)}
\label{sec:enzyme}

The founding analogy is salivary enzyme. It has \emph{no nucleus} and does not
replicate itself; it \emph{breaks down} food but does not build; it cannot
repair or rebuild itself and is replaced from outside, by the gland, at a
\emph{constant --- not exponential} rate; and it nonetheless does a job as a
system without being a conscious organism.

\begin{ruleblock}
\textbf{Design rule.} A verifying component \MUST{} be isolated from the thing
it verifies, and \MUSTNOT{} be able to repair itself. A verifier that can
rewrite itself can be made to agree with anything. Replacement of a verifier
comes from outside the verification path, at a bounded constant rate.
\end{ruleblock}

This is the structural argument for the three-channel separation of
\S\ref{sec:channels}.

\subsection{Filter and Flash}
\label{sec:filterflash}

\textbf{\textsc{filter}} --- sorting, searching, merging, splicing, slicing.
Filter operations are \emph{pure}: they mutate \emph{data} but \MUSTNOT{}
mutate \emph{program state}. Merge sort and splice are the canonical filters;
splice prunes an index without touching program state.

\textbf{\textsc{flash}} --- storing, sending, clearing, deleting. Flash is the
phase that \emph{commits}: it writes the message out, or writes it down.

Neither is meaningful alone. Filter without flash computes and discards; flash
without filter commits noise. The ratio is fixed by the \CISCO{} budget of
\S\ref{sec:ratio}: \textbf{two filter operations per one flash operation}.
Marshalling (\S\ref{sec:marshal}) is a flash operation --- \emph{``if you can
marshal, you're doing a flash.''}

% =====================================================================
\section{The Trident Packet}

\subsection{Why three}

The world is three-dimensional; therefore an interpretation served over it is
three-dimensional. A \textbf{trident packet} carries the message across three
axes --- $X$, $Y$, $Z$ --- and preserves it under any one of them. The three
axes are \emph{not} three copies; they are three \emph{readings} of one
message, which is why two of the three suffice (\S\ref{sec:threetwo}).

\subsection{Frame geometry}
\label{sec:frame}

The canonical frame derives from the reference string \OBI{}:

\begin{center}
\begin{tabular}{@{}llc@{}}
\toprule
\textbf{Quantity} & \textbf{Value} & \textbf{Derivation} \\
\midrule
Reference symbol count & 8 & \texttt{O B I N E X U S} \\
Index space per axis & 8 & indices $0$--$7$ \\
\textbf{Frame size} & \textbf{512 bytes} & $8^{3}$ \\
Bytes per symbol slot & 64 & $8^{2} = 512/8$ \\
Trident total & 1536 bytes & $512 \times 3$ \\
Half-frame & 256 bytes & $512/2$ \\
\textbf{Allocation quantum} & \textbf{4 bytes} & $256/64$ \\
\bottomrule
\end{tabular}
\end{center}

Memory \MUST{} be allocated ahead of time, in 4-byte sequences
(\S\ref{sec:budget}).

\begin{errata}{Errata A.2, A.3}
The generalisation to the full Latin alphabet is stated in S2 as $15{,}625 =
25^{3}$, using the \emph{maximum index} $25$ as the base. For 26 symbols
indexed $0$--$25$ the correct figure is $26^{3} = 17{,}576$. The 8-symbol case
$8^{3} = 512$ is correct as stated.
\end{errata}

\subsection{Axis pairs and sequence}

Two axes generate four ordered readings: $XX$, $XY$, $YX$, $YY$. $XX$ is the
\textbf{verification axis}; $XY$ and $YX$ carry direction of flow.

\begin{ruleblock}
\emph{``You cannot have a $Y$-axis without an $X$-axis.''} $Y$ is computable
from two $X$ readings (one positive, one negative) via the polar relation of
\S\ref{sec:lapis}. $Z$ is the projection --- the tangent relation between $X$
and $Y$ --- used for \emph{rotation}, not for payload. Therefore:
\textbf{verification operates on the $X$ axis only.}
\end{ruleblock}

\subsection{Index-pair addressing}
\label{sec:index}

Positions are addressed as \emph{index pairs}, not scalars:
\[
(0,0)\;(1,1)\;(2,2)\;(3,3)\;(4,4)\;(5,5)\;(6,6)\;(7,7)
\]
for the 8 slots of \OBI{}. The pair is \emph{(my index, your index)}: the
diagonal $(n,n)$ asserts sender/receiver agreement on slot $n$.

An off-diagonal pair received where a diagonal was expected is a \emph{detected
corruption}. Receiving $(3,2)$ where the schema says $(3,3)$ makes the message
wrong \emph{automatically}, with no inspection of the payload.

\textbf{Negative indices} are the polar inverse: $(-3,3)$ addresses the same
slot from the opposite end of the axis. When a receiver cannot resolve a symbol
at all, the sender transmits both poles so the receiver can locate the slot from
whichever end it reads from --- which is why a repair transmission carries up to
three packets, two at a time (\S\ref{sec:threetwo}).

\subsection{Token triplet}
\label{sec:triplet}

\begin{center}
\begin{tabular}{@{}ll@{}}
\toprule
\textbf{Field} & \textbf{Meaning} \\
\midrule
token type  & the type of the value \\
token value & the value itself \\
token data  & the datum the value acts on / is logged as \\
\bottomrule
\end{tabular}
\end{center}

Type and value are independently transmissible, which is what allows a receiver
to repair a value while keeping its type (\S\ref{sec:repair}).

\subsection{The $3 \rightarrow 2$ verification mapping}
\label{sec:threetwo}

A trident sends three messages; verification consumes two. With each labelled
\texttt{O} (order) or \texttt{C} (chaos), the receiver may be handed
$\texttt{OO}$, $\texttt{OC}$, $\texttt{CO}$, $\texttt{CC}$ --- four ordered
pairs, $\binom{3}{2} = 3$ unordered combinations.

\begin{ruleblock}
\textbf{Rule.} A receiver \MUST{} reach a verdict from any two of the three, and
\MUSTNOT{} require all three. The third is the spare --- it is what lets the
system survive a single-channel loss, and why the corpus notes that \emph{``two
homogeneous transistors can fail''} while the system still resolves.
\end{ruleblock}

The $3 \rightarrow 2$ relation is modelled on spring physics: the message is
pushed and returns. The operational consequence is the space/time duality of
\S\ref{sec:duality}.

\subsection{Corruption classes}
\label{sec:corrupt}

The corpus works a single example throughout:

\begin{Verbatim}[fontsize=\small,frame=single]
sent:      O    B    I    N    E    X    U    S
received:  O    I    e    t    e    X    U    S
                    (e = epsilon, t = theta)
\end{Verbatim}

\begin{center}
\begin{tabular}{@{}llp{6.0cm}l@{}}
\toprule
\textbf{Class} & \textbf{Symbol} & \textbf{Meaning} & \textbf{Operator} \\
\midrule
Substitution & wrong symbol & receiver resolved, resolved wrongly & \textbf{repair} \\
Empty & $\epsilon$ & receiver resolved to nothing & \textbf{renew} \\
Uncertain & $\theta$ & receiver unsure between candidates & \textbf{renew} + table \\
\bottomrule
\end{tabular}
\end{center}

$\epsilon$ is \emph{``you don't know what this is.''} $\theta$ is \emph{``you're
not sure what this is''} --- a \emph{bounded} ambiguity where the receiver has
candidates but cannot choose. The worked $\theta$ case: \texttt{N} may read as
\texttt{E}, \texttt{I} or \texttt{M} depending on where the reader started on
the lookup grid, because \emph{``my N is your E, my E is your X, my X is your
S''} --- a uniform off-by-one shift of the reader's frame.

\subsection{Repair procedure}
\label{sec:repair}

Repair is \textbf{index-first, payload-second}. The sender \MUSTNOT{}
retransmit the payload for a slot the receiver resolved correctly.

\begin{enumerate}[leftmargin=*]
  \item \texttt{O} at slot 0 --- correct. \textbf{No action.} \emph{``O is fine.
    I verified. I didn't need to send O again.''}
  \item \texttt{B} at slot 1 received as \texttt{I} --- substitution. Transmit
    the \emph{index pair} $(1,1)$, not the letter.
  \item $\epsilon$ at slot 2 --- empty. \textbf{Renew}: transmit $(2,2)$ plus
    the value.
  \item $\theta$ at slot 3 --- uncertain. Transmit the polar pair $(-3,3)$ plus
    a shift-of-view instruction (\S\ref{sec:shift}).
  \item \texttt{X}, \texttt{U}, \texttt{S} --- correct. \textbf{No action.}
\end{enumerate}

\begin{ruleblock}
\textbf{Schema, not message.} The verification target is the \emph{schema}, not
the payload: \emph{``the schema is supposed to verify, not the message\dots{}
message can be tampered with but not the schema.''} This is why index pairs are
the repair currency --- they are schema-level assertions.
\end{ruleblock}

\subsection{Lookup grid and shift-of-view}
\label{sec:shift}

Where $\theta$ cannot be resolved by index alone, the sender transmits a
\textbf{shift-of-view}: a relative move on a 2-D lookup grid of the symbol set,
as $(\Delta\text{row}, \Delta\text{col})$. To move from \texttt{I} to \texttt{N}
is \emph{``one right, one down''}; the inverse is $(-1,-1)$. The sender \MAY{}
transmit the single move, the pair, or --- as a last resort --- the whole
substructure of the table.

Permitted grid transformations are the standard four --- \emph{rotation,
reflection, translation, enlargement} --- parameterised by $\theta$, with a full
rotation $2\pi$. The search space is the area swept; \S\ref{sec:parity} bounds
it.

\subsection{XOR parity}
\label{sec:xor}

Verification across the trident uses \textbf{XOR}, on the ground that XOR
\emph{remembers state}: it encodes the prior value into the result and is its
own inverse. The corpus works the hex directly:
\[
\texttt{O} = \texttt{0x4F}, \qquad
\texttt{B} = \texttt{0x42}, \qquad
\texttt{I} = \texttt{0x49}.
\]
Where a slot survives unchanged, $v \oplus v = 0$ and the parity contributes
nothing; where it differs, the parity carries the difference. \emph{``When you
verify a message over XOR you get another message --- a parity message --- and
the parity sequence gives you the result you need.''}

Polar sign convention: a right-hand-side XOR preserves the byte; the negative
operator is the left-hand-side reading of the same structure. The poles
\emph{multiply} rather than add, because they are dimensional rather than
scalar.

% =====================================================================
\section{\CISCO{} Order}
\label{sec:cisco}

\subsection{Definition}

\CISCO{} denotes the corpus's bottom-up self-balancing tree discipline. It has
nothing to do with the network vendor. A \CISCO{} tree is: self-balancing, and
specifically a red-black tree \emph{that never needs pruning}; traversed
$\text{root} \rightarrow \text{left} \rightarrow \text{right}$; parsed
\emph{bottom-up} while the top-down model shares the same two poles at each
step; and is the structure over which \OBI{} is held, with \texttt{O} at the
root, \texttt{B} left, \texttt{I} right.

The Riemann hypothesis is invoked in S3 as the intended means of deciding, from
signal structure alone, whether a number was produced by bottom-up traversal.
\emph{(Unresolved --- Appendix~\ref{app:open}, C.1.)}

\subsection{Read / Write / Execute permutations}

Three permissions generate six orderings:

\begin{Verbatim}[fontsize=\small,frame=single]
R W X        W R X        X R W
R X W        W X R        X W R
\end{Verbatim}

The \emph{diagonals} are where a permission is shared between two positions ---
\texttt{R R} on one, \texttt{W W} on the other. Those shared cells make the
matrices composable: two matrices sharing a diagonal can be verified against
each other. Octal permission tags follow POSIX and are carried on the wire:
$r = 4$, $w = 2$, $x = 1$.

\begin{errata}{Errata A.10}
S5b states ``one would be easy to write.'' Under POSIX, $1$ is \emph{execute}
and $2$ is \emph{write}.
\end{errata}

\subsection{The $2R \rightarrow 1W \rightarrow 1X$ ratio}
\label{sec:ratio}

\begin{align*}
2 \;\text{reads} &\;\rightarrow\; 1 \;\text{write} \\
2 \;\text{writes} &\;\rightarrow\; 1 \;\text{execute}
\end{align*}

One read alone is a write in progress; \emph{two} reads complete a write. A
write is thus one \emph{implicit} filter plus one \emph{explicit} flash
(\S\ref{sec:filterflash}). In frame terms: $512$ bytes per $2$ read operations
$=$ $1$ execute operation.

\subsection{Hops, and rotation instead of traversal}

Moving between R, W and X costs \emph{hops}: minimum one, maximum two left hops
to the execute root.

\begin{ruleblock}
\textbf{The two hops are not serial and are not to be walked.} The tree is
\emph{rotated}: \emph{``you don't have to do the hops per se, you just have to
do one context switch\dots{} you just change the angle of the message.''} Two
hops become two rotations of the \CISCO{} tree; the message is unchanged, only
the reader's frame moves.
\end{ruleblock}

Under rotation: my execute $\rightarrow$ your read; my read $\rightarrow$ your
write; my write $\rightarrow$ your execute. The dimension is preserved because
the flash was already sent --- only the angle changed. This is the same
mechanism as \S\ref{sec:shift}, applied to the permission tree rather than the
symbol grid.

\subsection{Determinant sharing}
\label{sec:det}

Two $2\times3$ (or $3\times2$) permission matrices are reconciled by exchanging
\emph{transposes}, not the matrices themselves, and computing a determinant
relation over the pair. The worked figures: $R_w = [1,0,1]$, $R_r = [0,1,2]$,
with transposes exchanged so the shared cells ($0$ and $1$) are visible to both
parties without either disclosing its full matrix.

\begin{errata}{Errata A.11}
The determinant arithmetic in S2 is not recoverable --- the audio degrades into
repeated self-correction and the stated products do not reconcile. The
\emph{mechanism} (exchange transposes, verify via a determinant relation,
disclose no full matrix) is clear and is what is normative. The specific figures
are not.
\end{errata}

% =====================================================================
\section{Channel Model}
\label{sec:channels}

\subsection{Roles}

\begin{center}
\begin{tabular}{@{}llp{7.6cm}@{}}
\toprule
\textbf{Ch} & \textbf{Role} & \textbf{Function} \\
\midrule
CH0 & Transmitter & encodes raw message into packet format; computes SHA-256
  message hash; attaches r/w/x permission tag and order/chaos status \\
CH1 & Receiver & deserialises the packet; decodes from the wire form \\
CH2 & Verifier & brokers consensus; broadcasts the verdict \\
\bottomrule
\end{tabular}
\end{center}

The three roles map onto the only three verbs the system has ---
\emph{receive, transmit, verify} --- in all six orderings. Channels \MUST{} be
independently verifiable and \MAY{} be run independently. The isolation
principle (\S\ref{sec:enzyme}) requires that CH2 cannot repair itself.

\begin{errata}{Inconsistency}
S5b describes the ring as ``channel 0 sender, channel 2 receiver, channel 3
verifier'' in one passage and demonstrates \texttt{make ch0} / \texttt{make ch1}
with a channel-2 verifier in another. This document normalises on
CH0/CH1/CH2.
\end{errata}

\subsection{Ring topology and worked exchange}

\begin{Verbatim}[fontsize=\small,frame=single]
        CH1 (receiver)
       /               \
   CH0 ----------------- CH2
 (sender)            (verifier)

CH0 ->  "Hello. I am here for food."
CH1 ->  "Yes. Here is your food."
CH2 ->  "Verified."
\end{Verbatim}

\subsection{Observed failure modes}
\label{sec:failures}

The S5b demonstration produced three failures. All are \emph{specified
behaviour} and \MUST{} be reported rather than masked.

\begin{center}
\begin{tabular}{@{}lp{10.2cm}@{}}
\toprule
\textbf{Failure} & \textbf{Meaning} \\
\midrule
\texttt{hash mismatch} & SHA-256 digest of received frame $\neq$ transmitted
  digest. The message may still be legible; the frame is not verified. \\
\texttt{consensus failed} & CH2 could not reach a verdict from two of three.
  The system \MUSTNOT{} exit; it holds. \\
\texttt{phantom ID} & An identifier appeared on the ring with no originating
  channel. Not fatal --- observed, not acted on. \\
\bottomrule
\end{tabular}
\end{center}

In the demonstration the message \emph{arrived and was legible} while the hash
failed. That is the intended separation: legibility and verification are
different verdicts.

\subsection{RGB tomographic state}
\label{sec:rgb}

Channel state is carried as an RGB triple whose components are \emph{pointers to
state}, not colours:

\begin{center}
\begin{tabular}{@{}lp{10.0cm}@{}}
\toprule
\textbf{Red} & \textbf{drift.} A drifting party appears more red. Two reds
  (light/dark) are distinguished by hue, saturation, lightness. \\
\textbf{Green} & \textbf{verified.} Green is the consistency check. A message is
  \emph{not green until it is sent and verified}. \\
\textbf{Blue} & \textbf{neutral / character state.} The observing state ---
  neither drifted nor verified. \\
\bottomrule
\end{tabular}
\end{center}

The three parties share the \emph{dimension} (hue), not the value. Agreement is
the qualification between two reds; green is the anchor state measurable by a
third party. The corpus frames this explicitly as \textbf{zero-knowledge}: the
verifier confirms agreement without holding either party's value. Grey/brown is
the decayed state --- information loss, measured by ratio over the shared
$\pi$-frequency.

\subsection{Consensus and channel eviction}

Where a channel becomes hostile or is captured, the remaining channels
\emph{vote it out}. The worked case: \texttt{A4}, \texttt{B4}, \texttt{C4}
observe that \texttt{D3} lied to \texttt{B4}; the three vote, \texttt{D3} is
crossed out, and connections are killed until the real-world cause is
established. Eviction \MUST{} be by consensus, never unilateral; \MUST{} be
reversible on re-establishment of cause; and \SHOULD{} use XOR over the open
channel (\S\ref{sec:xor}).

\subsection{Actors}

\textbf{Obi} --- heart, soul; king or leader. The \emph{local observer}: sees
the system from the inside out, observes tomographically, and is \emph{not
directly involved in the command and control}. \textbf{Uche} --- knowledge. The
\emph{acting party}. Each can create and destroy the other. Neither is the
verifier of itself.

% =====================================================================
\section{Signal Model}

\subsection{The noise/signal quadrants}
\label{sec:quadrants}

\begin{center}
\begin{tabular}{@{}lll@{}}
\toprule
 & \textbf{signal} & \textbf{no signal} \\
\midrule
\textbf{noise}    & noise $+$ signal    & noise, no signal \\
\textbf{no noise} & no noise $+$ signal & \textbf{no noise, no signal} \\
\bottomrule
\end{tabular}
\end{center}

\begin{itemize}[leftmargin=*,noitemsep]
  \item \textbf{no noise, no signal} --- the system is on, potential,
    quiescent. \emph{Not} off.
  \item \textbf{no noise $+$ signal} --- clean channel; connection may be
    established.
  \item \textbf{noise $+$ signal} --- live and degraded; the working case.
  \item \textbf{noise, no signal} --- losing connection, or drifting out of
    detection.
\end{itemize}

Machine states map on: \emph{on}, \emph{standby}, \emph{hibernate} (regulating
resources), \emph{sleep} (off for a short time).

\subsection{Your noise is my signal}

\begin{ruleblock}
\emph{``Your noise is my signal.''} An adversary jamming a band is emitting
structure, and that structure is information about the adversary. A receiver
\SHOULD{} treat detected jamming as a signal source, not only as degradation.
\end{ruleblock}

Consequently, \textbf{a complete system is one that knows it can be jammed}:
\emph{``a system that can be jammed knows it can be jammed --- so it can be
anti-jammed.''} Self-knowledge of the failure mode is a completeness criterion.

\subsection{The maybe lattice}

The two binaries of \S\ref{sec:quadrants} cross with the
\emph{maybe} / \emph{maybe-not} operator to give a $4 \times 4$ state lattice.
\textbf{maybe-not} is a \emph{nullifying} operator: it nullifies a state in
order to verify it --- ``check for noise; test for no noise.'' It is the
interrogative form of the state, which makes the lattice a calibration procedure
rather than a classification.

\emph{Balloon model.} Each channel is a balloon that inflates and deflates.
Inflation is signal emerging from noise; deflation is signal collapsing into
noise. Three balloons --- maybe, maybe-not, and the popped/ground case --- give
the trident its physical reading.

\begin{errata}{Errata A.9}
The worked figure in S4b, $2^{-1} \times 3 = -0.1$, does not reconcile:
$2^{-1}\times 3 = 1.5$ and $(-2)^{-1}\times 3 = -1.5$. The intended quantity is
not recoverable.
\end{errata}

\subsection{\textsc{ox-star} and \textsc{noise-star}}
\label{sec:oxstar}

\textbf{\textsc{ox-star}} (auxiliary star) is the on-the-fly injection
mechanism: a sequence that turns a system on, off, or into a new mode
\emph{while it is running}, without a human at the location. The canonical
sequence is the Konami code:
\[
\uparrow\;\uparrow\;\downarrow\;\downarrow\;\leftarrow\;\rightarrow\;\leftarrow\;\rightarrow\;\texttt{B}\;\texttt{A}\;\texttt{START}
\]

This is not decoration. The corpus's reading is precise and is the protocol's
own send-rule:

\begin{itemize}[leftmargin=*,noitemsep]
  \item $\uparrow$ once is \emph{half} the instruction --- held, not sent.
  \item $\uparrow\uparrow$ is the complete instruction --- \textbf{sent}.
    Likewise $\downarrow\downarrow$, $\leftarrow\rightarrow$.
  \item \texttt{B A} are the bounded terminals --- a stop-and-place.
  \item \texttt{START} commits; \texttt{SELECT} requests more options.
\end{itemize}

This is the \emph{half the structure, double the send} rule
(\S\ref{sec:duality}) made concrete: an \textsc{ox-star} is a scalar tensor that
takes half a dimension up, so each instruction is pressed twice to be sent once.

\textbf{\textsc{noise-star}} is the complementary state family ---
\texttt{ox-noise}, \texttt{ox-no-noise}, \texttt{ox-signal},
\texttt{ox-no-signal} --- carrying the \S\ref{sec:quadrants} quadrant into the
injection layer, so an injected command knows the standby state of the system it
is entering. Stated applications: guiding a drone in real time; disarming a
system mid-flight; switching a running program's mode without unloading it.

\subsection{Probing}
\label{sec:probe}

\textbf{Processing} is executing what the system already knows how to do.
\textbf{Probing} is the system seeking a \emph{new} solution structure for the
same problem. A probe \MUSTNOT{} be logged as a process.

\begin{align*}
\text{external probe:}&\quad P(\text{state}) \rightarrow \text{data} \\
\text{internal probe:}&\quad P(\text{data}) \rightarrow \text{state}
\end{align*}

A probe asks the five interrogatives --- \emph{who, what, when, where, why} ---
and does so \textbf{without a human in the loop}. The corpus distinguishes
\emph{human in the loop}, \emph{human on the loop} and \emph{human out of the
loop}; probing is out-of-loop by construction.

Probe health is a scalar in $[0,1]$ on each node. A structurally healthy node
probes for more inputs; a degraded node probes for more statements. This makes
protocol functions \emph{dynamic} --- able to resolve ambiguity by
interpretation. The minimal case --- \emph{``I am Python.'' ``No, I am
Python.'' Who is right?} --- is the ambiguity a probe exists to settle.

% =====================================================================
\section{Here-and-Now, There-and-Then, Where-and-Whenever}
\label{sec:addressing}

\subsection{The temporal-spatial matrix}

The protocol addresses place and time as a $2 \times 3$ matrix over three modes
(S4a, S5; triad ratified by the author 25 August 2026). The canonical statement:

\begin{ruleblock}
\begin{center}
\textbf{here and now \quad--\quad there and then \quad--\quad where and whenever}
\end{center}
\end{ruleblock}

Each mode names a \emph{degree of binding}, and each has two orderings --- space
first, or time first --- which is what makes the matrix $2 \times 3$ rather than
a list of three.

\begin{center}
\begin{tabular}{@{}>{\raggedright\arraybackslash}p{2.5cm}>{\raggedright\arraybackslash}p{2.0cm}>{\raggedright\arraybackslash}p{2.2cm}>{\raggedright\arraybackslash}p{3.6cm}>{\raggedright\arraybackslash}p{3.6cm}@{}}
\toprule
\textbf{Mode} & \textbf{Binding} & \textbf{Invariant} & \textbf{space first} & \textbf{time first} \\
\midrule
\textsc{here and now} & both bound, \emph{deictic} & never a toy &
  \emph{here and now} --- present in space, then time &
  \emph{now and here} --- present in time, then space \\[2pt]
\textsc{there and then} & both bound, \emph{displaced} & never a weapon &
  \emph{there and then} --- there in space, then in time &
  \emph{then and there} --- there in time, then in space \\[2pt]
\textsc{where and whenever} & both \emph{unbound} & never a problem &
  \emph{where and whenever} --- in space, for all time &
  \emph{whenever and where} --- in time, for all space \\
\bottomrule
\end{tabular}
\end{center}

Each mode carries exactly one of the Three Nevers (\S\ref{sec:nevers}). The
invariants are not a separate governance layer --- they are the conformance
conditions on the addressing modes.

Each cell is a \emph{dimensionless quantity} occurring on \emph{one axis only}
--- one place and one time at a time. The corpus likens the arrangement to a
gimbal, and to latitude/longitude, and states both orderings of the first mode
explicitly: \emph{``I need food here and now, or now and here.''}

\subsection{The binding progression}

The three modes are not three locations. They are three \emph{quantifier states}
over the same pair of dimensions.

\begin{enumerate}[leftmargin=*]
  \item \textsc{here and now} --- place $=$ \emph{this}, time $=$ \emph{this}.
    Both anchored to the speaker. The corpus gives this mode a \textbf{loopback
    address}, so a system can address its own present.
  \item \textsc{there and then} --- place $=$ \emph{that}, time $=$ \emph{that}.
    Displaced, but still \emph{specified}: a particular elsewhere, a particular
    past or future.
  \item \textsc{where and whenever} --- place and time are \emph{quantified},
    not specified. This is the \textbf{interrogative} mode: \emph{where} and
    \emph{when} are interrogative pronouns, and \emph{whenever} is their
    free-choice form. The claim is open, and resolves at whatever place and time
    the need arises.
\end{enumerate}

The third mode is where \textbf{probing} lives (\S\ref{sec:probe}). A probe asks
\emph{who, what, when, where, why} --- and the two dimensions it quantifies over
are exactly the two this mode leaves unbound.

\subsection{Stillness}
\label{sec:stillness}

\textbf{Stillness} is not a fourth mode. It is the \emph{property} the unbound
mode has: the state in which a claim persists without further action ---
\emph{``you don't need anything to be still.''} The demonstration is the cup: a
drink set down out of frame, still there when you return, returnable to
indefinitely.

\begin{ruleblock}
\textbf{Normative consequence.} A claim entered under \textsc{where and
whenever} \MUSTNOT{} expire, and \MUSTNOT{} be discharged by the passing of a
specific place or time. \emph{``I need food at this address at 18:00 Tuesday''}
is a bound claim and can be missed. \emph{``I need food''} --- unbound --- is a
\textbf{standing claim}: honoured wherever the claimant is and whenever the need
arises. Food ordered under stillness arrives when it is needed, not when it was
asked for.
\end{ruleblock}

This is what makes a duty-holder unable to discharge an obligation by pointing
at a missed slot, and it is precisely the content of this mode's invariant,
\emph{\textbf{never a problem}} (\S\ref{sec:nevers}): an unbound claim cannot
time out into nothing, because it has no slot to miss and therefore accrues no
debt. It is not failing; it is unresolved.

Its one failure mode is \textbf{relaxation} --- a standing claim that quietly
stops being honoured --- which is why relaxation is specified as a \emph{failure}
and not as an \emph{expiry}. Relaxation is the only way the unbound mode can
become a problem, and is therefore the only thing this invariant has to guard.

\subsection{Structural note}

Three modes $\times$ two orderings $=$ \textbf{six cells}, matching the $3+3$
dimensional count (three axes forward, three back) and the six R/W/X
permutations of \S\ref{sec:cisco}. Whether these three sixes are the same six,
or a coincidence of structure, is left open
(Appendix~\ref{app:open}, C.8).

\subsection{Parity}
\label{sec:parity}

Parity is a \emph{dimension}, not a predicate:
\[
x \bmod 2 = 0 \;\Rightarrow\; \text{even}, \qquad
x \bmod 2 = 1 \;\Rightarrow\; \text{odd}.
\]
A parity check runs on \emph{one axis} and determines the even/odd ratio for
that axis. A negative quantity ($-2$) is not merely a sign: the minus is an
\emph{operator on the negative axis}, under which ordinary parity rules bend.
Encoding \MUST{} therefore be checkable for bit-flip pairs on both poles. Where
the rules bend, the transformation set of \S\ref{sec:shift} applies --- and the
edge cases where the rules break are precisely what the stateless protocol must
encode, so that statelessness survives them.

\subsection{The holding problem}
\label{sec:holding}

A system that computes an intermediate result \emph{holds it in memory}, and
holding is the cost. Worked as a matrix product: the intermediate outer-product
table is what the system is forced to hold, when what it wants is the row sums.

\begin{ruleblock}
\textbf{Goal.} Look up the result from a tabulation constructed by polar rules,
rather than compute-and-hold.
\end{ruleblock}

\begin{errata}{Errata A.6}
The intermediate table stated in S6 --- rows $(6,10)$, $(12,15)$, $(16,20)$
with row sums $16, 27, 36$ --- is \emph{not} a consistent outer product of any
single vector pair. For $u = (2,3,4)$ and $v = (3,5)$ the correct outer product
is rows $(6,10)$, $(9,15)$, $(12,20)$ with row sums $16, 24, 32$. The
\emph{intent} --- memoise intermediates so the result is a lookup, not a
computation --- is sound and is what is normative.
\end{errata}

\textbf{Gravity well / rope and bucket.} A well has a depth (density); a rope
holds a bucket; the rope \MUST{} hold indefinitely. \emph{``Nothing likes to
hold anything.''} Indefinite holding must therefore be \emph{structural}, not
effortful --- encoded into the geometry, not maintained by a process.

\subsection{\LAPIS{} --- polar quadrants}
\label{sec:lapis}

\LAPIS{} is the corpus's polar-calculus mnemonic for the four cardinal bearings.
As given in S2: $\pi/4 = $ north, $\pi/3 = $ east, $\pi/2 = $ south, $\pi = $
west.

\begin{errata}{Errata A.5}
This mapping is not consistent with standard polar convention, and is not
internally consistent either --- $\pi/4$ and $\pi/3$ are $45^\circ$ and
$60^\circ$, which are not cardinal. Under the convention the corpus itself
invokes (``the sun rises in the east and sets in the west''), the cardinal
bearings are
\[
0 = \text{east}, \quad \pi/2 = \text{north}, \quad
\pi = \text{west}, \quad 3\pi/2 = \text{south}.
\]
The \emph{intent} --- four quadrants addressed by angle, with a message choosing
its bearing $\theta$ before it moves --- is preserved. Only the constants are
wrong. Implementations \MUST{} use the standard convention.
\end{errata}

The associated rule is real and independent of the constants: \emph{determine
the angle you want to move before you move}, and $\theta$ (rotation) and
$\lambda$ (power distribution) are chosen together.

\subsection{The $i$-cycle and the trident mismatch}
\label{sec:icycle}

S6 develops a rotation table over the imaginary unit:
\[
i^{1} = i, \quad i^{2} = -1, \quad i^{3} = -i, \quad i^{4} = 1
\qquad \text{(period 4)}
\]

\begin{errata}{Errata A.7 --- and a proposed correction}
A period-\emph{4} cycle is a poor algebraic fit for a \emph{trident}. The
natural object for a three-axis rotation is the cube roots of unity:
\[
\omega = e^{2\pi i/3}, \qquad \omega^{3} = 1, \qquad 1 + \omega + \omega^{2} = 0.
\]
The identity $1 + \omega + \omega^{2} = 0$ is exactly the property a trident
wants: three readings summing to nothing when the message is intact, and to a
non-zero residue when it is not. That is a parity check for free, on three axes
--- which is what \S\ref{sec:xor} reaches for with XOR. Offered as an
improvement, not as a transcription of the corpus.
\end{errata}

\subsection{Dimensional principality}

\begin{ruleblock}
\textbf{Axiom.} Dimensions have principles they obey without being told. A
system \MUSTNOT{} need to be given time to think in order to respect them ---
the rules are implicit and hold structurally.
\end{ruleblock}

Non-linearity follows: \emph{any structure that can be computed by looking at it
must be isomorphic under transposition}. The transpose is an \emph{axis of
reflection} --- $XYZ \rightarrow ZYX$ --- under which a row vector becomes a
column vector and the data structure is preserved while the system moves.

The $3+3$ count: three axes forward ($x,y,z$) and three back ($-x,-y,-z$),
because what is behind you is real whether or not you can see it. The corpus
demonstrates this with a mirror: the cup out of frame is still there.

\begin{errata}{Errata A.8}
S6 states $3 + 6 + 9 = 15$ and derives a ``minimum consensus'' from it:
$3+6+9 = 18$. It also states $3^{3} = 9$: in fact $3^{2} = 9$ and $3^{3} = 27$.
$6 \times 9 = 54$ is correct as stated.
\end{errata}

% =====================================================================
\section{Codec and File Formats}
\label{sec:codec}

\subsection{\texttt{.nsigii} --- symbolic expression of intent}

The protocol's file format is a \emph{symbolic expression of intent} file:

\begin{itemize}[leftmargin=*,noitemsep]
  \item It holds space \emph{indefinitely} --- once a point is allocated the
    program holds it. It cannot be removed, only \emph{reallocated}.
  \item It is \emph{stateless but remembering}. Statelessness here does not mean
    ``does not remember''; it means the state is encoded into the \emph{geometry}
    of the system rather than held in a variable. \emph{``It remembers stuff, it
    doesn't need to forget, ever.''}
  \item New symbol sets can be added as language evolves, without invalidating
    existing files.
\end{itemize}

\subsection{Linkable-then-executable coding}

The \NSIGII{} codec is \textbf{LTE-coding --- linkable \emph{then} executable}.
The demonstrated pipeline:

\begin{Verbatim}[fontsize=\small,frame=single]
cat image.png | go run main.go   ->  raw, linkable, stateless encoding
\end{Verbatim}

The codec takes \emph{any} input format and produces a raw encoding that
\emph{``can never die''} --- it survives re-encoding. The output is
\emph{linkable as a library} and \emph{executable on demand}: a DLL, or a
\emph{macro system for other systems}, not a direct call. It encodes the
\emph{buffering}, which is the point --- a conventional stream buffers before it
plays; \NSIGII{} encodes the buffer so the receiver already holds the state.

Stated repositories: \texttt{github.com/obinexus/nuko-operating-system}
(\textsc{nuko-os}, the human-rights operating system) and
\texttt{github.com/obinexus/legislation}.

\subsection{Marshalling}
\label{sec:marshal}

Marshalling is the codec's serialisation contract, explicitly compared to
Python's \texttt{pickle}:

\begin{ruleblock}
A receiver holding the \emph{codec of the function} can reconstitute an object
it did not construct, because it knows the \emph{shape} of the function without
computing it.
\end{ruleblock}

Marshalling is a \textbf{flash} operation (\S\ref{sec:filterflash}). Comparable
wire formats named in the corpus: Protocol Buffers, gRPC, JSON, and a binary IR.
\NSIGII{} is schema-based like these, which is what \S\ref{sec:repair}'s
schema-verification requires.

% =====================================================================
\section{Dimensional Classifiers --- the Natural Language Layer}
\label{sec:lang}

\subsection{Statement / Expression / Intent}

Three poles, likened to latitude, longitude and altitude:

\begin{center}
\begin{tabular}{@{}lll@{}}
\toprule
\textbf{Pole} & \textbf{Carries} & \textbf{Truth-apt?} \\
\midrule
Statement  & the assertion                   & yes --- true or false \\
Expression & the meaning of the assertion    & no --- it means \\
Intent     & the abstract model behind it    & no --- it aims \\
\bottomrule
\end{tabular}
\end{center}

The bipolar question of an expression is \emph{``what could this mean?''} --- and
that is where a natural-language codec must operate. The worked pair
\emph{John likes Mary} / \emph{Mary likes John} are two statements over one
relation, whose conjunction is positive or negative by pole. The corpus reaches
for the discriminant $b^{2} - 4ac$ as the object deciding how many resolutions
such a relation has. \emph{(Unresolved --- Appendix~\ref{app:open}, C.3.)}

\subsection{Verb--noun construction}
\label{sec:verbnoun}

Following Igbo construction, the instruction unit is a \textbf{verb--noun pair}.
A \emph{verb} describes an action (``to'' $+$ action); an \emph{adverb}
describes the verb and situates it in past, present or future; a \emph{noun} is
the classifier of the thing --- and a \emph{person} is only one category of
noun. The corpus is deliberate here: nouns include animals and things, so intent
can be read from non-human experience.

\subsection{The four sentence types}

\begin{center}
\begin{tabular}{@{}llp{5.6cm}@{}}
\toprule
\textbf{Type} & \textbf{Example} & \textbf{Protocol role} \\
\midrule
Declarative   & \emph{I am Nnamdi.} & assertion --- carries a statement \\
Interrogative & \emph{Can you introduce yourself?} & \textbf{probe} (\S\ref{sec:probe}) --- requests input \\
Imperative    & \emph{Close the door.} & \textbf{command} --- carries control \\
Exclamative   & \emph{What a funny thing he told us!} & unexpected-state marker \\
\bottomrule
\end{tabular}
\end{center}

Crossed with four sentence structures: short/simple; simple (personal);
compound (one dependent $+$ one independent clause); complex (two independent
$+$ one dependent); compound-complex. The corpus draws a \textsc{cisc}/\textsc{risc}
analogy across the classifiers: the reduced and complex sets operate on
different instruction models, but a receiver must resolve intent from either.

\subsection{Clauses as control flow, and the \texttt{then} marker}
\label{sec:then}

\begin{Verbatim}[fontsize=\small,frame=single]
do  ->  while           do { x++ } while (x <= 25)
do  ->  then  ->  while do { x++ } then LABEL while (...)
if  ->  then  ->  else  ->  then  ->  if
\end{Verbatim}

\begin{ruleblock}
The \texttt{then} keyword \textbf{labels code blocks}, and labelled blocks are
what a \texttt{goto} targets --- so \texttt{then} is a \textbf{switchable goto
marker}.
\end{ruleblock}

\texttt{then} is part of the language construct. It \MUSTNOT{} be switched out
or broken by a consuming implementation; a language that cannot express it
cannot interpret \NSIGII{} clauses. The stated implementation target is
\textbf{gosilang}.

Traversal order matters and is carried: pre-order (\texttt{++x}) and post-order
(\texttt{x++}) produce different ranges over the same loop, and the difference
is an \emph{execution-order} property of the clause, not an idiom.

% =====================================================================
\section{Sizing and Timing Model}

\subsection{Frame budget}
\label{sec:budget}

\begin{center}
\begin{tabular}{@{}ll@{}}
\toprule
\textbf{Quantity} & \textbf{Value} \\
\midrule
Frame & 512 bytes ($8^{3}$) \\
Trident & 1536 bytes ($512 \times 3$) \\
Half-frame & 256 bytes \\
Pre-allocation & 2048 bytes ($512 \times 4$) \\
Allocation quantum & 4 bytes \\
Digest & 32 bytes (SHA-256) \\
Payload after digest & 480 bytes \\
\bottomrule
\end{tabular}
\end{center}

Memory \MUST{} be allocated ahead of time, in 4-byte sequences --- \emph{per}
four bytes, not \emph{over} four bytes. The allocation is a grid, and the grid
is what is searched for meaning.

\begin{errata}{Errata A.1 --- the trident split}
$512/3 = 170.\overline{6}$, not an integer: the trident split of a 512-byte
frame across three axes does \emph{not} divide evenly. Implementations \MUST{}
choose one of:
\begin{itemize}[leftmargin=*,noitemsep]
  \item a \textbf{510-byte payload} split $170/170/170$, with 2 bytes reserved
    for parity --- preserving the corpus's $170 = \texttt{0b10101010}$
    alternating pattern, which is itself load-bearing (\S\ref{sec:index});
    \textbf{or}
  \item an uneven split $171/171/170$.
\end{itemize}
The first is recommended: $170 = \texttt{10101010}_2$ is the alternating index
pattern the addressing model already uses, and the 2-byte remainder gives the
XOR parity of \S\ref{sec:xor} somewhere to live.
\end{errata}

\subsection{Space--time duality}
\label{sec:duality}

\[
\text{half the space} \;\longleftrightarrow\; \text{double the time},
\qquad
\text{double the space} \;\longleftrightarrow\; \text{half the time}
\]

A message may be halved and sent over twice the interval, or doubled and sent in
half. \textbf{The worst case is the time case} --- the corpus optimises for
bounded latency, not bounded memory. The spring model gives the mechanism: a
message is \emph{pushed} and \emph{returns}, and the return is what verifies it.
This is why the $3\rightarrow2$ mapping works: the third packet is the spring's
return path.

\subsection{Bit-rate derivation}

\begin{errata}{Errata A.4 --- corrected derivation}
S2 states ``2 seconds $=$ 6,000 milliseconds'' and derives
$6000/1024 = 5.859375 \rightarrow 5.86$. The division is arithmetically correct;
the \emph{input} is not. $2$ seconds $=$ \textbf{2,000\,ms}, not 6,000 --- the
transcript applies a seconds-to-minutes factor of 60 where none belongs.
Corrected, for a 512-byte frame over a 2-second window:
\begin{align*}
\text{throughput} &= 512\,\text{B} / 2\,\text{s} = 256\,\text{B/s} = 2{,}048\,\text{bit/s} \\
\text{doubled frame} &= 1{,}024\,\text{B} \\
\text{per-unit interval} &= 2{,}000\,\text{ms} / 1{,}024 = 1.953125\,\text{ms} \\
\text{two-unit interval} &= 3.90625\,\text{ms}
\end{align*}
S2's later figure of $51.72$\,ms derives from $25.86 \times 2$, where $5.86$ was
transcribed as $25.86$. Neither figure survives the corrected input.
\end{errata}

\subsection{Spring and damping}

Springs model \emph{storing and absorbing} energy --- a message pushed against
resistance stores what it will need to return.

\begin{center}
\begin{tabular}{@{}>{\raggedright\arraybackslash}p{6.4cm}>{\raggedright\arraybackslash}p{7.6cm}@{}}
\toprule
\textbf{Spring failure} & \textbf{Protocol analogue} \\
\midrule
fatigue --- repeated cycles grow microscopic cracks & accumulated retransmission degrades the channel \\
overload --- permanent deformation or breakage & a frame beyond budget does not recover \\
corrosion --- environmental degradation over time & schema drift \\
relaxation --- loss of tension over time & a stillness claim that stops being honoured \\
\bottomrule
\end{tabular}
\end{center}

\begin{errata}{Errata A.12 (unresolved)}
A damping factor of $2.85$ is stated in S3, attributed to ``Roth theory,
clustering, nearest neighbour.'' A damping \emph{ratio} $\zeta = 2.85$ is
heavily overdamped ($\zeta = 1$ is critical). Neither ``Roth theory'' nor the
relation of $2.85$ to the clustering/nearest-neighbour method is recoverable
from the transcript. Flagged for the author.
\end{errata}

% =====================================================================
\section{Governance and Enforcement Layer}

\subsection{The pointer problem}

\begin{ruleblock}
\textbf{Breathing and living are pointers that \MUST{} hold in all contexts.
Work is optional.}
\end{ruleblock}

A pointer that must not be null is the thing a system is built to guarantee. If
the breathing pointer drifts, someone is dying. If the living pointer drifts,
someone is not living. Work is a pointer that \MAY{} be null --- and a system
that forces work by dereferencing it anyway \emph{``causes fatigue and
confusion.''}

\NSIGII{} is therefore a \textbf{human-rights broker}, not a work broker. It
holds the demographic state, and it holds \emph{negative space for work} ---
because work becomes possible when space and time are already held.

\subsection{Litigation as enforcement}
\label{sec:litigation}

S5b develops the enforcement mechanism concretely: \NSIGII{} composes a
verified, timestamped, hash-anchored record of a request for food, water or
shelter, and of the duty-holder's response or non-response. The record is
transmitted by ordinary means --- the corpus demonstrates it over SMTP --- and
is intended to be usable as evidence that a statutory duty was engaged and not
discharged.

Statutes named as the intended encoding targets (UK): Care Act 2014 (and the
Health and Social Care Act 2012); Children Act 1989, ss.~17, 21, 23C
(accommodation; care leavers); Housing Act; Equality Act 2010. The stated design
goal is that a person can assert an entitlement \emph{without first knowing the
statute}, because the protocol encodes which duty a given request engages.

\begin{errata}{Note}
This specification describes an evidentiary and record-keeping mechanism.
Whether a given record supports a given claim is a question of law and of the
specific facts, and needs a qualified solicitor --- a protocol cannot supply
that, and this document does not.
\end{errata}

\subsection{Reconciliation of ``never a weapon''}
\label{sec:reconcile}

The mode pairing of \S\ref{sec:nevers} resolves the tension by giving it a
locus. ``Never a weapon'' is the invariant of \textsc{there and then} --- the
displaced mode, where all of the protocol's reach across space and time lives.

\begin{ruleblock}
\textbf{``Never a weapon'' scopes to the protocol's own reach.} \NSIGII{}
\MUSTNOT{} be constructible into an instrument that acts at a distance to harm
--- it cannot be turned against its user, cannot surveil them, cannot deny them.
Its \textsc{ox-star} injection, remote command and mid-flight control
capabilities (\S\ref{sec:oxstar}) are all there-and-then operations and all fall
under this prohibition.

It does \emph{not} scope to the actions a claimant takes with the record the
protocol produces. A person using their own verified record to enforce their own
statutory entitlement is acting in \textsc{here and now}, under the invariant
\emph{never a toy} --- a requirement that the matter be treated as serious, not
a prohibition on acting. That is the protocol working, not the protocol
weaponised.
\end{ruleblock}

The two are separable because they are invariants of \emph{different modes}. The
protocol's reach is constrained; the claimant's presence is not.

\emph{Ratified in structure by the author's mode/never pairing (25 August 2026);
the wording above remains open to refinement.}

% =====================================================================
\section{Conformance}

An implementation conforms to \NSIGII{} 0.3 if it:

\begin{enumerate}[leftmargin=*,noitemsep]
  \item Implements all six operators of \S\ref{sec:ops}, or declares which half
    of a pair is structurally impossible and why.
  \item Carries an \texttt{ORDER}/\texttt{CHAOS} status on every channel and
    treats \texttt{CHAOS} as an observable rather than a fault.
  \item Isolates the verifier such that it cannot repair itself.
  \item Maintains the 2-filter-to-1-flash ratio.
  \item Emits trident packets of the frame geometry in \S\ref{sec:frame} and
    \S\ref{sec:budget}, with the $512/3$ split resolved per Errata A.1.
  \item Addresses slots by index pair and repairs index-first.
  \item Distinguishes substitution, $\epsilon$ and $\theta$, applying the
    matching operator to each.
  \item Reaches a verdict from any two of three readings and never requires all
    three.
  \item Verifies the \emph{schema}, not the payload.
  \item Implements CH0/CH1/CH2 as independently runnable, and reports
    \texttt{hash mismatch}, \texttt{consensus failed} and
    \texttt{phantom ID} rather than masking them.
  \item Holds --- does not exit --- on consensus failure.
  \item Requires consensus of remaining channels for eviction.
  \item Never mistakes \emph{no noise, no signal} for \emph{off}.
  \item Logs probes distinctly from processes.
  \item Addresses every claim in one of the three modes of
    \S\ref{sec:addressing}, and \textbf{never expires a claim entered under}
    \textsc{where and whenever} (\S\ref{sec:stillness}).
  \item Allocates memory ahead of time in 4-byte quanta.
  \item Satisfies the Three Nevers of \S\ref{sec:nevers} \textbf{as mode
    invariants} --- \emph{never a toy} on \textsc{here and now}, \emph{never a
    weapon} on \textsc{there and then}, \emph{never a problem} on \textsc{where
    and whenever} --- and does not treat them as a separate governance layer.
\end{enumerate}

% =====================================================================
\appendix
\section{Arithmetic Errata}
\label{app:errata}

Every numeric claim in the corpus, checked. \emph{Stated} is what the transcript
says; \emph{Correct} is the computed value; the disposition preserves the
intent.

{\small
\begin{longtable}{@{}>{\raggedright\arraybackslash}p{0.6cm}>{\raggedright\arraybackslash}p{2.6cm}>{\raggedright\arraybackslash}p{2.9cm}>{\raggedright\arraybackslash}p{2.9cm}>{\raggedright\arraybackslash}p{5.1cm}@{}}
\toprule
\textbf{\#} & \textbf{Claim} & \textbf{Stated} & \textbf{Correct} & \textbf{Disposition} \\
\midrule
\endhead
A.1 & Trident split of a frame & $512/3$ treated as clean & $170.\overline{6}$ --- does not divide & 510-byte payload $170/170/170$ $+$ 2 parity bytes; $170 = \texttt{0b10101010}$ preserved \\
A.2 & Frame size & $8^{3} = 512$ & $512$ \checkmark & Correct as stated \\
A.3 & Latin index space & $15{,}625 = 25^{3}$ & $26^{3} = 17{,}576$ & Off-by-one: max index 25 used as base instead of symbol count 26 \\
A.4 & Milliseconds in 2\,s & $6{,}000$\,ms & $2{,}000$\,ms & Seconds$\rightarrow$minutes factor of 60 wrongly applied \\
A.4b & Derived interval & $5.86 \rightarrow 51.72$\,ms & $1.953125$\,ms; $\times 2 = 3.90625$\,ms & $51.72$ comes from $25.86\times2$; $5.86$ mis-transcribed as $25.86$ \\
A.5 & \LAPIS{} bearings & $\pi/4$\,N, $\pi/3$\,E, $\pi/2$\,S, $\pi$\,W & $0$\,E, $\pi/2$\,N, $\pi$\,W, $3\pi/2$\,S & Not internally consistent; intent (bearing chosen before movement) preserved \\
A.6 & Outer-product table & rows $(6,10)$, $(12,15)$, $(16,20)$; sums $16$, $27$, $36$ & for $(2,3,4)\otimes(3,5)$: rows $(6,10)$, $(9,15)$, $(12,20)$; sums $16$, $24$, $32$ & Stated table is not a consistent outer product; memoisation intent sound \\
A.7 & Rotation cycle & $i,-1,-i,1$ (period 4) & correct for $i$, but a trident wants period 3 & Use $\omega = e^{2\pi i/3}$; $1+\omega+\omega^{2}=0$ gives a 3-axis parity check \\
A.8 & Consensus minimum & $3+6+9=15$; $3^{3}=9$ & $3+6+9=18$; $3^{2}=9$, $3^{3}=27$; $6\times9=54$ \checkmark & Two slips; the 54 is right \\
A.9 & Probe scalar & $2^{-1}\times3 = -0.1$ & $1.5$; or $-1.5$ for $(-2)^{-1}$ & Not recoverable --- flagged \\
A.10 & POSIX octals & ``1 would be easy to write'' & $r{=}4$, $w{=}2$, $x{=}1$ & $1$ is execute, $2$ is write \\
A.11 & Determinant exchange & figures do not reconcile & --- & Mechanism normative; figures are not \\
A.12 & Damping factor & $2.85$, ``Roth theory'' & --- & $\zeta = 2.85$ heavily overdamped; attribution unrecoverable \\
A.13 & Suffering equation & $\Sigma = (N-R)\times K$ & needs clamping, finite $K$ & As stated, $R \geq N$ gives $\Sigma \leq 0$; $N = R$ with $K\to\infty$ is $0\times\infty$ \\
A.14 & Suffering, corrected & --- & $\Sigma = K\cdot\max(0, N-R)$, $K \in (0,\infty)$ & Preserves both stated behaviours. Recommended form \\
A.15 & Happiness model & $H = S + C + \dots$ & $H = S + C + V$ & Lyubomirsky--Sheldon--Schkade: set point $S$, circumstances $C$, voluntary activity $V$; original weighting contested \\
A.16 & Trident total & $512\times3$ & $1{,}536$ \checkmark & Correct \\
A.17 & Pre-allocation & $512\times4$ & $2{,}048$ \checkmark & Correct \\
A.18 & Half-frame, quantum & $512/2$; $256/64$ & $256$; $4$ \checkmark & Correct \\
A.19 & $170$ in binary & \texttt{10101010} & $\texttt{0b10101010} = 170$ \checkmark & Correct; the alternating pattern is load-bearing \\
A.20 & Hex of symbols & \texttt{4F}, \texttt{42}, \texttt{49} & \texttt{0x4F}, \texttt{0x42}, \texttt{0x49} \checkmark & Correct ASCII \\
A.21 & Konami sequence & stated correctly & $\uparrow\uparrow\downarrow\downarrow\leftarrow\rightarrow\leftarrow\rightarrow$ \texttt{B A} \checkmark & Correct \\
A.22 & Digest & ``SH 56'' & SHA-256, 32-byte digest & Leaves 480 payload bytes in a 512-byte frame \\
\bottomrule
\end{longtable}
}

\subsection*{A.14 --- the suffering equation, in full}

The corpus (S4a) defines experienced suffering over need $N$, resources $R$ and
a constraint constant $K$, with two required behaviours: suffering collapses
toward zero when resources meet need, and grows without bound as the constraint
grows. As stated, $\Sigma = (N-R)K$ satisfies neither cleanly. The corrected
form:
\[
\Sigma \;=\; K \cdot \max\bigl(0,\; N - R\bigr), \qquad K \in (0, \infty)
\]
\begin{itemize}[leftmargin=*,noitemsep]
  \item $R \geq N \;\Rightarrow\; \Sigma = 0$ --- no negative suffering.
  \item $N > R$, $K$ large $\;\Rightarrow\; \Sigma$ large --- constraint
    dominates, as intended.
  \item $N = R$ with $K \to \infty$ is now $0$, not the undefined
    $0 \times \infty$.
\end{itemize}

The companion happiness model stated alongside it is
$H = S + C + V$ --- genetic set point, circumstances, voluntary activity ---
which is the Lyubomirsky--Sheldon--Schkade formulation. Its original 50/10/40
weighting is contested in the later literature and \SHOULD{} not be treated as
a protocol constant.

% =====================================================================
\section{Reference Alphabet}

\begin{center}
\begin{tabular}{@{}lll@{}}
\toprule
\textbf{Letter} & \textbf{Standard NATO} & \textbf{Corpus usage} \\
\midrule
O & Oscar & Oscar \\
B & Bravo & Bravo \\
I & India & India / \emph{Indigo} \\
N & November & November \\
E & Echo & Echo \\
X & X-ray & X-ray \\
U & Uniform & \emph{Umbrella} \\
S & Sierra & \emph{Sarah} \\
G & Golf & Golf \\
C & Charlie & Charlie \\
H & Hotel & Hotel \\
\bottomrule
\end{tabular}
\end{center}

Reference string \OBI{} --- 8 symbols, indices $0$--$7$, the basis of the
512-byte frame. Protocol name \NSIGII{} --- November, Sarah/Sierra,
Indigo/India, Golf, India, India.

% =====================================================================
\section{Open Questions}
\label{app:open}

\begin{description}[leftmargin=1.6cm,style=nextline]
\item[C.1] \textbf{Riemann and bottom-up decidability.} S3 asserts the Riemann
  hypothesis ``allows you to verify whether the number was determined via
  \CISCO{}, which is bottom-up.'' No mechanism connects the distribution of zeta
  zeros to traversal order. Either a concrete construction is needed, or the
  claim should be restated as an analogy.
\item[C.2] \textbf{Damping factor 2.85.} Source and role unrecovered.
\item[C.3] \textbf{Discriminant in the language layer.} $b^{2}-4ac$ is invoked
  to decide the resolution count of a two-party relation. The mapping from a
  declarative pair to $a$, $b$, $c$ is not given.
\item[C.4] \textbf{``Never a weapon'' vs.\ litigation. \textsc{resolved} in
  v0.3.} The mode pairing scopes ``never a weapon'' to \textsc{there and then}
  --- the protocol's own reach --- while a claimant enforcing their own record
  acts in \textsc{here and now} under \emph{never a toy}. The two are invariants
  of different modes and are therefore separable. Wording in
  \S\ref{sec:reconcile} remains open to refinement.
\item[C.5] \textbf{Channel numbering.} CH2 vs CH3 as verifier --- normalised to
  CH2, pending confirmation.
\item[C.6] \textbf{Igbo tone in the wire format.} The tonal \emph{eze} example
  is the motivating case for schema-level verification, but no encoding for tone
  is specified. If tone is meaning, the frame needs a tone channel --- this may
  be what the $Z$ axis is for.
\item[C.7] \textbf{Cube roots of unity.} Whether to adopt $\omega$ in place of
  $i$ is a design decision, not a correction. Recommended, but the author's
  call.
\item[C.8] \textbf{The three sixes.} Three addressing modes $\times$ two
  orderings, the $3+3$ dimensional count, and the six R/W/X permutations all
  give six. Whether these are the same six under different names --- which would
  let a claim's addressing mode be derived from its permission tag, or vice
  versa --- or a coincidence of structure, is open. If they are the same six,
  the addressing mode need not be carried on the wire at all.
\end{description}

% =====================================================================
\section{Traceability}

\begin{center}
\begin{tabular}{@{}ll@{}}
\toprule
\textbf{Spec section} & \textbf{Corpus segment} \\
\midrule
\S2.1 Name, Igbo tone & S3, S5b \\
\S2.3 Three Nevers (mode invariants) & S5; pairing ratified 25 Aug 2026 \\
\S2.5 Toolchain & S3, S6b, S5b \\
\S3.1 Six operators & S1 \\
\S3.2 Order/Chaos & S1 \\
\S3.3 Enzyme model & S1 \\
\S3.4 Filter/Flash & S2, S5b \\
\S4 Trident packet & S1, S2 \\
\S4.7 Corruption classes & S1 \\
\S4.10 XOR parity & S1, S4a \\
\S5 \CISCO{} order & S1, S2 \\
\S6 Channel model & S5b \\
\S6.4 RGB tomographic & S3 \\
\S6.5 Eviction & S4b \\
\S6.6 Obi / Uche & S3 \\
\S7.1--7.2 Signal model & S3, S4b \\
\S7.3 Maybe lattice & S4b \\
\S7.4 \textsc{ox-star} & S3 \\
\S7.5 Probing & S3, S4a \\
\S8.1--8.3 Addressing triad & S4a, S5; ratified 25 Aug 2026 \\
\S8.3 Holding problem & S6a \\
\S8.4 \LAPIS{} & S2 \\
\S8.5--8.6 Dimensionality & S6a \\
\S9 Codec & S6b, S5b \\
\S10 Language layer & S6b \\
\S11 Sizing and timing & S1, S2, S3 \\
\S12 Governance & S3, S5, S5b \\
\bottomrule
\end{tabular}
\end{center}

\vfill
\begin{center}
\emph{End of \NSIGII{} Protocol Specification v0.3 (Draft).}
\end{center}

\end{document}
